%! AP Statistics — Unit 8, Lesson 8.4 — Guided Notes
\documentclass[10pt]{article}
\usepackage{apstats-article}
\usepackage{apstats-boxes}
\newcommand{\TallMath}[1]{$\displaystyle #1\rule[-1.4em]{0pt}{3.2em}$}

\begin{document}

\pageheader{Unit 8, Lesson 8.4}{Guided Notes}
\namedateperiod

\begin{objectivebox}
By the end of this lesson, I will be able to:
\begin{itemize}
    \item[$\square$] \writeline
    \item[$\square$] \writeline
    \item[$\square$] \writeline
\end{itemize}
\end{objectivebox}

\begin{vocabbox}
\termblanklong{Two-way table}
\termblanklong{Expected count (two-way)}
\termblanklong{Marginal distribution}
\end{vocabbox}

\begin{hookbox}
\textbf{Discussion:} A survey asks 200 students about their grade level and learning
preference (online vs.\ in-person). Without doing any calculation, do you think grade
level and learning preference are associated? Why or why not?

\writelines{2}
\end{hookbox}

\begin{notesbox}{Reading a Two-Way Table}
In a two-way table:
\begin{itemize}
    \item Rows represent categories of the variable \blank{5cm}
    \item Columns represent categories of the variable \blank{5cm}
    \item Each cell contains the \blank{4cm} of individuals in that row-and-column combination
    \item The \textbf{row totals} and \textbf{column totals} form the \blank{5cm} distribution
\end{itemize}

\renewcommand{\arraystretch}{1.6}
\begin{tabular}{lcccc}
             & \textbf{Online} & \textbf{In-person} & \textbf{Row Total} \\
    \hline
    9th grade  & 28 & 22 & \blank{1.5cm} \\
    10th grade & 35 & 15 & \blank{1.5cm} \\
    11th grade & 22 & 28 & \blank{1.5cm} \\
    12th grade & 15 & 35 & \blank{1.5cm} \\
    \hline
    Column Total & \blank{1.5cm} & \blank{1.5cm} & \blank{1.5cm} \\
\end{tabular}
\end{notesbox}

\begin{notesbox}{What Are Expected Counts?}
Under $H_0$ (independence or homogeneity), knowing someone's row category tells us
\blank{6cm} about their column category.

\vspace{0.1in}
\textbf{The Expected Count Formula:}
\[
    \text{Expected count} = \frac{(\text{row total}) \times (\text{column total})}{\blank{3cm}}
\]

\textbf{Derivation --- two equivalent routes:}

\renewcommand{\arraystretch}{2.0}
\begin{tabular}{ll}
    \textbf{Route A} (row proportion $\times$ column total): &
        $\dfrac{\text{row total}}{\blank{3cm}} \times \text{column total}$ \\
    \textbf{Route B} (row total $\times$ column proportion): &
        $\text{row total} \times \dfrac{\text{column total}}{\blank{3cm}}$ \\
\end{tabular}

\vspace{0.1in}
Both simplify to $\dfrac{(\text{row total})(\text{column total})}{\text{table total}}$.
\end{notesbox}

\begin{notesbox}{Computing Expected Counts: Grade/Learning Example}
Use the table above. For the 9th grade / Online cell:

\vspace{0.1in}
$E = \dfrac{(\blank{1.5cm}) \times (\blank{1.5cm})}{\blank{1.5cm}} = \blank{2cm}$

\vspace{0.1in}
Fill in all 8 expected counts:

\renewcommand{\arraystretch}{1.8}
\begin{tabular}{lccc}
             & \textbf{Online} & \textbf{In-person} & \textbf{Row Total} \\
    \hline
    9th grade  & \blank{2cm} & \blank{2cm} & 50 \\
    10th grade & \blank{2cm} & \blank{2cm} & 50 \\
    11th grade & \blank{2cm} & \blank{2cm} & 50 \\
    12th grade & \blank{2cm} & \blank{2cm} & 50 \\
    \hline
    Col Total  & 100 & 100 & 200 \\
\end{tabular}

\vspace{0.1in}
\textbf{Large-Counts Condition:} All expected counts must be $\ge$ \blank{1cm}.

Check: Is the condition met here? \blank{5cm}
\end{notesbox}

\begin{practicebox}
\textbf{Guided Practice:} A random sample of 90 college students is classified by year
(First-year / Sophomore / Junior) and whether they live on campus (Yes / No).

\renewcommand{\arraystretch}{1.6}
\begin{tabular}{lccc}
             & \textbf{On Campus} & \textbf{Off Campus} & \textbf{Row Total} \\
    \hline
    First-year  & 22 & 8  & 30 \\
    Sophomore   & 15 & 15 & 30 \\
    Junior      & 8  & 22 & 30 \\
    \hline
    Col Total   & 45 & 45 & 90 \\
\end{tabular}

\vspace{0.1in}
Compute the expected count for each of the 6 cells. Show work for at least 2 cells.

\begin{enumerate}
  \item First-year / On Campus: $E = \dfrac{(\blank{1.5cm})(\blank{1.5cm})}{\blank{1.5cm}} = \blank{2cm}$
        \vspace{0.05in}
  \item Sophomore / Off Campus: $E = \blank{6cm} = \blank{2cm}$
        \vspace{0.05in}
  \item Fill in the remaining 4 cells using the same formula.
\end{enumerate}

Check the large-counts condition: \writeline
\end{practicebox}

\end{document}
